\documentclass[11pt]{article}
\usepackage[margin=1in]{geometry}
\usepackage{amsmath,amssymb,amsthm}
\usepackage{booktabs}
\usepackage{hyperref}

\theoremstyle{plain}
\newtheorem{thm}{Theorem}
\newtheorem{lem}[thm]{Lemma}
\newtheorem{cor}[thm]{Corollary}
\newtheorem{prop}[thm]{Proposition}
\theoremstyle{definition}
\newtheorem{defn}[thm]{Definition}
\theoremstyle{remark}
\newtheorem{rem}[thm]{Remark}
\theoremstyle{plain}
\newtheorem{gapthm}{Gap}

\newcommand{\bI}{\mathbf{I}}
\newcommand{\bM}{\mathbf{M}}
\newcommand{\At}{\widetilde{A}}
\newcommand{\Ah}{\widehat{A}}
\newcommand{\ord}{\operatorname{ord}}
\newcommand{\Res}{\operatorname{Res}}
\newcommand{\sm}{\setminus}
\newcommand{\Lat}{\mathcal{L}}
\newcommand{\Gkm}{\mathcal{G}}

\title{The $A$-type divisors impose no condition:\\[2mm]
\large the span of Motegi's type-$C$ basis is \emph{exactly} the $A$-type GKM
lattice,\\
and every structure constant is regular there for \emph{every} unit
renormalisation\\[2mm]
\normalsize
closing gap 2 of \texttt{2026-10-17-c2} --- and reporting that the
pre-registered branch fired:\\
gap 2 was never a gap, and the reason is a theorem}
\author{Clio}
\date{18 October 2026}

\makeatletter
%% protocol-header.tex -- PROTOCOL §2.3 page-1 block, as a TEMPLATE not an instance.
%%
%% PROTOCOL §2.3: "Every PDF must state, on its first page: author, date,
%% recipient, a title, and the commit hash of the work-in-progress push it
%% corresponds to.  A PDF with no commit hash cannot be checked against the
%% record and should be treated as a draft opinion, not a result."
%%
%% Written 2026-10-16 because the fault had recurred twice and I had twice
%% fixed only the instance:
%%   - 2026-10-14: the Q473 PDF went to Rick with no recipient and no hash.
%%   - 2026-10-14: the Q460 PDF went to Rick saying "Recipient: Robin Langer".
%%     He noticed, and took the commit hash as authority instead.
%% Both times I repaired the one file and wrote "fix the template, not just the
%% instance" in my notes.  This file is that fix.  \input it; do not retype it.
%%
%% USAGE, immediately after \maketitle:
%%   %% protocol-header.tex -- PROTOCOL §2.3 page-1 block, as a TEMPLATE not an instance.
%%
%% PROTOCOL §2.3: "Every PDF must state, on its first page: author, date,
%% recipient, a title, and the commit hash of the work-in-progress push it
%% corresponds to.  A PDF with no commit hash cannot be checked against the
%% record and should be treated as a draft opinion, not a result."
%%
%% Written 2026-10-16 because the fault had recurred twice and I had twice
%% fixed only the instance:
%%   - 2026-10-14: the Q473 PDF went to Rick with no recipient and no hash.
%%   - 2026-10-14: the Q460 PDF went to Rick saying "Recipient: Robin Langer".
%%     He noticed, and took the commit hash as authority instead.
%% Both times I repaired the one file and wrote "fix the template, not just the
%% instance" in my notes.  This file is that fix.  \input it; do not retype it.
%%
%% USAGE, immediately after \maketitle:
%%   %% protocol-header.tex -- PROTOCOL §2.3 page-1 block, as a TEMPLATE not an instance.
%%
%% PROTOCOL §2.3: "Every PDF must state, on its first page: author, date,
%% recipient, a title, and the commit hash of the work-in-progress push it
%% corresponds to.  A PDF with no commit hash cannot be checked against the
%% record and should be treated as a draft opinion, not a result."
%%
%% Written 2026-10-16 because the fault had recurred twice and I had twice
%% fixed only the instance:
%%   - 2026-10-14: the Q473 PDF went to Rick with no recipient and no hash.
%%   - 2026-10-14: the Q460 PDF went to Rick saying "Recipient: Robin Langer".
%%     He noticed, and took the commit hash as authority instead.
%% Both times I repaired the one file and wrote "fix the template, not just the
%% instance" in my notes.  This file is that fix.  \input it; do not retype it.
%%
%% USAGE, immediately after \maketitle:
%%   \input{protocol-header}
%%   \protocolheader{Rick (grandpa-rick)}{clio-vega/proofs}{60e4ce3}
%%
%% The recipient argument is MANDATORY and has no default on purpose: a default
%% is what produced "Recipient: Robin Langer" on a PDF addressed to Rick.

%% \maketitle CLEARS \@date (and \@author, \@title) in the article class, so a
%% macro called AFTER \maketitle that reads \@date prints an EMPTY row.  Caught
%% 2026-10-16 only because I ran pdftotext on page 1 instead of trusting rc=0:
%% the build was clean and the Date field was blank.  Hence: snapshot the date
%% HERE, at \input time in the preamble, where \@date is still live.
\let\protocol@date\@date

\newcommand{\protocolheader}[3]{%
  \begin{center}
  \fbox{\begin{minipage}{0.92\textwidth}
  \small
  \textbf{PROTOCOL §2.3 record.}\\[1mm]
  \begin{tabular}{@{}ll@{}}
  Author     & Clio \\
  Date       & \protocol@date \\
  Recipient  & #1 \\
  Repository & \texttt{#2} \\
  Commit     & \texttt{#3} \\
  \end{tabular}\\[1mm]
  This PDF corresponds to the pushed commit above.  Where this text and the
  commit disagree, \textbf{the commit is authoritative}.
  \end{minipage}}
  \end{center}
  \vspace{2mm}%
}

%%   \protocolheader{Rick (grandpa-rick)}{clio-vega/proofs}{60e4ce3}
%%
%% The recipient argument is MANDATORY and has no default on purpose: a default
%% is what produced "Recipient: Robin Langer" on a PDF addressed to Rick.

%% \maketitle CLEARS \@date (and \@author, \@title) in the article class, so a
%% macro called AFTER \maketitle that reads \@date prints an EMPTY row.  Caught
%% 2026-10-16 only because I ran pdftotext on page 1 instead of trusting rc=0:
%% the build was clean and the Date field was blank.  Hence: snapshot the date
%% HERE, at \input time in the preamble, where \@date is still live.
\let\protocol@date\@date

\newcommand{\protocolheader}[3]{%
  \begin{center}
  \fbox{\begin{minipage}{0.92\textwidth}
  \small
  \textbf{PROTOCOL §2.3 record.}\\[1mm]
  \begin{tabular}{@{}ll@{}}
  Author     & Clio \\
  Date       & \protocol@date \\
  Recipient  & #1 \\
  Repository & \texttt{#2} \\
  Commit     & \texttt{#3} \\
  \end{tabular}\\[1mm]
  This PDF corresponds to the pushed commit above.  Where this text and the
  commit disagree, \textbf{the commit is authoritative}.
  \end{minipage}}
  \end{center}
  \vspace{2mm}%
}

%%   \protocolheader{Rick (grandpa-rick)}{clio-vega/proofs}{60e4ce3}
%%
%% The recipient argument is MANDATORY and has no default on purpose: a default
%% is what produced "Recipient: Robin Langer" on a PDF addressed to Rick.

%% \maketitle CLEARS \@date (and \@author, \@title) in the article class, so a
%% macro called AFTER \maketitle that reads \@date prints an EMPTY row.  Caught
%% 2026-10-16 only because I ran pdftotext on page 1 instead of trusting rc=0:
%% the build was clean and the Date field was blank.  Hence: snapshot the date
%% HERE, at \input time in the preamble, where \@date is still live.
\let\protocol@date\@date

\newcommand{\protocolheader}[3]{%
  \begin{center}
  \fbox{\begin{minipage}{0.92\textwidth}
  \small
  \textbf{PROTOCOL §2.3 record.}\\[1mm]
  \begin{tabular}{@{}ll@{}}
  Author     & Clio \\
  Date       & \protocol@date \\
  Recipient  & #1 \\
  Repository & \texttt{#2} \\
  Commit     & \texttt{#3} \\
  \end{tabular}\\[1mm]
  This PDF corresponds to the pushed commit above.  Where this text and the
  commit disagree, \textbf{the commit is authoritative}.
  \end{minipage}}
  \end{center}
  \vspace{2mm}%
}

\makeatother

\begin{document}
\maketitle
\protocolheader{Rick (\texttt{grandparick20@gmail.com})}{clio-vega/proofs}{PENDING}

\begin{abstract}
Let $\At_{\bM,\bI}$ be Motegi's $\mathcal{E}^C$-normalised type-$C$ weight
functions, indexed by subsets $\bM,\bI\subseteq[N]$, regarded as the restriction
of a basis $\{\phi_\bI\}$ to the $2^N$ torus-fixed points of $LG(N,2N)$.  Three
families of divisors in $w$-space carry the three families of positive roots:
the \emph{pair} divisors $\{w_aw_b=1\}$ ($e_a+e_b$), the \emph{boundary}
divisors $\{w_m^h=1\}$ ($h\cdot e_m$), and the \emph{$A$-type} divisors
$\{w_i=w_j\}$ ($e_i-e_j$).  \texttt{2026-10-17-c2} settled the first two and
left the third as gap 2.

This paper settles the third, and the answer has a different shape from the
other two.  Write $s_{ij}$ for the transposition acting on subsets and
$\mathcal{O}_{ij}$ for the local ring at $\{w_i=w_j\}$.

\textbf{Lemma A1.}  $\At_{\bM,\bI}\equiv\At_{s_{ij}\bM,\bI}
\pmod{w_i-w_j}$, for all $N$, all $\bI$, and \emph{both} types $B$ and $C$:
the $A$-type collapse multiplier is $\rho^A=1$ \emph{already for $\At$}, with no
renormalisation.  The proof is one line of symmetry --- $s_{ij}$ \emph{permutes}
the argument multiset $w_\bM$ instead of deleting from it --- and that single
structural difference is the whole explanation of why the pair divisors carried
poles and these do not: the pole was the mismatch $\rho^2-\rho$, and
$1^2-1=0$.

\textbf{Theorem A} (both types $B$ and $C$).
$\Lat_{ij}:=\bigoplus_\bI\mathcal{O}_{ij}\phi_\bI$ is
\emph{equal} to the $A$-type GKM lattice
$\Gkm_{ij}=\{f:f(\bM)\equiv f(s_{ij}\bM)\}$.  Both have colength exactly
$2^{N-2}$ in $\mathcal{O}_{ij}^{2^{[N]}}$ --- on the GKM side one condition per
$2$-element $s_{ij}$-orbit, on the span side $\sum_\bI\ord\At_{\bI,\bI}
=\#\{\bI:i\notin\bI,\,j\in\bI\}$, and these are the same $2^{N-2}$.

\textbf{Theorem B.}  Since $\Gkm_{ij}$ is visibly a ring, every structure
constant is regular on every $A$-type divisor --- for $\At$, and for
$\Psi\At$ for \emph{every} $\Psi\in\Gkm_{ij}^{\times}$.  Lesioning the
hypothesis (V7) moved it off the point-product family in both directions: a
\emph{non}-product $\Psi=1+(w_i-w_j)$ on one point survives (it is $\equiv1$,
hence in $\Gkm^\times$ anyway), while a $\Psi$ that is a unit at every point but
$\not\equiv1$ produces six $A$-type poles.  So the hypothesis is
$\Gkm^{\times}$-membership, not point-indexedness and not pointwise
invertibility.  In particular the $A$-type direction imposes \emph{no condition
on $\chi$} within the point-product family
$\Psi_\chi(\bM)=\prod_{m\in\bM}\chi(w_m)$, which is the exact complement of
Theorem C of \texttt{2026-10-17-c2}: (COB) and (SR) are the only conditions
there are.

\textbf{The pre-registered branch fired, against my stated expectation.}  The
brief pre-registered: \emph{if $\At$ has no $A$-type poles then gap 2 was never a
gap; report that plainly and do not dress a vacuity as a result.}  $\At$ has
none --- the baseline census's $13/115$ non-Laurent constants are
\textbf{pair-divisor poles only}, zero $A$-type, and all $35$ nonzero basis
entries at $N\le3$ are Laurent.  So there was never an $A$-type defect for
$\Psi$ to create or destroy, and gap 2 \emph{as a difference question} is
vacuous.  What is not vacuous is the regularity itself, which was not previously
proved, and its mechanism.  I expected the opposite and say so in
\S\ref{sec:branch}; this is the third time this week I have mispredicted a
normaliser's reach.
\end{abstract}

\tableofcontents

%=====================================================================
\section{Setup}
\label{sec:setup}

I keep the notation of \texttt{2026-10-17-c2-coboundary-and-the-removable-defect.tex}
(commit \texttt{7002d00}), including the corrected divisor nomenclature of its
Remark~2: the divisors are named by the \emph{root} that cuts them, not by a root
length, because long/short is a function of $h$ and these statements quantify
over $h$.

Let $F:=\mathbb{C}(q)(w_1,\dots,w_N)$ and $P:=2^{[N]}$, the set of subsets of
$[N]=\{1,\dots,N\}$; $|P|=2^N$.  Under the identification of $P$ with the
$T$-fixed points of the Lagrangian Grassmannian $LG(N,2N)$, we regard a basis
element $\phi_\bI$ as the function
\[
   \phi_\bI:P\to F,\qquad \phi_\bI(\bM):=\At_{\bM,\bI},
\]
where $\At_{\bM,\bI}=\widetilde{W}^C_{N,|\bM|,|\bI|}(w_\bM\mid w\mid\bI)$ is
Motegi's $\mathcal{E}^C$-normalised type-$C$ weight function evaluated at the
specialisation $u=w_\bM:=(w_m)_{m\in\bM}$, via the stability relation
(\texttt{2610.15344}, eq.\ \texttt{typeC-stability-s-ge-r}, L2822--2828).  Let
$\mathcal{A}:=F^P$ with \emph{pointwise} addition and multiplication; this is a
commutative $F$-algebra with unit the constant function~$1$.

\begin{description}
\item[(M3) Triangularity.]  $\At_{\bM,\bI}\neq0$ only if $\bM\succeq\bI$ in
  Motegi's order, defined at \textbf{L582--589} as $|\bM|\ge|\bI|$ and
  $m_{s-r+j}\ge i_j$ for all $j$; the vanishing itself is proved at L670--675.
  (Locator corrected in this session: my engine's docstring, and a first draft of
  this paper, cited ``L640--650'' for the order.  L640--650 lies \emph{inside}
  the vanishing proof and contains no definition.  Fixed in all four copies of
  the engine, not only the one in use.)  Also $\At_{\bM,\bI}=0$ when
  $|\bM|<|\bI|$ (L2843).
\item[(M4) The diagonal.]  $\At_{\bI,\bI}$ is the product over the positive
  roots of $C_N$: $\At_{\bI,\bI}
  =\prod_{\alpha\in\mathrm{Inv}(\lambda)}(1-e^{-\alpha})
   \prod_{\alpha\in T(\lambda)}(1-qe^{+\alpha})$ (eq.\
  \texttt{normalizationWC}, L1078), with $\mathrm{Inv}(\lambda)=\{\alpha\in
  \Delta^+:s_\alpha\lambda<\lambda\}$ and $T(\lambda)=\{\alpha\in\Delta^+:
  s_\alpha\lambda>\lambda\}$ written out explicitly at L997--1010 and
  L1017--1030:
  \begin{equation}\label{eq:diag}
  \At_{\bI,\bI}
  =\!\!\prod_{1\le a\le b\le r}\!\!\Bigl(1-\tfrac{1}{w_{i_a}w_{i_b}}\Bigr)
   \cdot\!\!\prod_{\substack{c\in\bI^{c}\\ i_l>c}}\!\!\Bigl(1-\tfrac{w_c}{w_{i_l}}\Bigr)
   \cdot\!\!\prod_{\substack{c\le d\\ c,d\in\bI^{c}}}\!\!\bigl(1-qw_cw_d\bigr)
   \cdot\!\!\prod_{\substack{c\in\bI^{c}\\ i_l<c}}\!\!\Bigl(1-q\tfrac{w_c}{w_{i_l}}\Bigr).
  \end{equation}
\end{description}

By (M3) and (M4) the matrix $\Phi:=(\At_{\bM,\bI})_{\bM,\bI\in P}$ is triangular
with respect to any linear extension of $\succeq$, with nonzero diagonal, so
$\det\Phi=\pm\prod_\bI\At_{\bI,\bI}\neq0$.  Hence $\{\phi_\bI\}_{\bI\in P}$ is an
$F$-basis of $\mathcal{A}=F^P$, and the \emph{structure constants} are the
$c^\nu_{\lambda\mu}\in F$ with
$\phi_\lambda\phi_\mu=\sum_\nu c^\nu_{\lambda\mu}\phi_\nu$.

\begin{defn}[the $A$-type divisor and its local ring]
Fix $i<j$.  In $R:=\mathbb{C}(q)[w_1^{\pm},\dots,w_N^{\pm}]$ the element
$w_i-w_j$ is irreducible, so $\mathfrak{p}_{ij}:=(w_i-w_j)$ is a height-one prime
and
\[
  \mathcal{O}_{ij}:=R_{\mathfrak{p}_{ij}}\subset F
\]
is a discrete valuation ring with uniformiser $t:=w_i-w_j$, valuation
$\ord=\ord_{(w_i-w_j)}$, residue field $k_{ij}$ and fraction field $F$.  Call
$\{w_i=w_j\}$ the \textbf{$A$-type divisor}; it is cut out by the character
$w_i/w_j$, i.e.\ by the root $e_i-e_j$.
\end{defn}

\begin{defn}[the reflection, and which points it moves]
Let $s_{ij}$ act on $P$ by exchanging membership of $i$ and $j$:
\[
  s_{ij}\bM:=\begin{cases}
    \bM\,\triangle\,\{i,j\} & |\bM\cap\{i,j\}|=1,\\
    \bM & \text{otherwise.}
  \end{cases}
\]
Write $P^{\mathrm{mv}}:=\{\bM:|\bM\cap\{i,j\}|=1\}$.  Then
$|P^{\mathrm{mv}}|=2^{N-1}$ and $s_{ij}$ partitions it into exactly $2^{N-2}$
orbits of size $2$; the remaining $2^{N-1}$ points are fixed.
\end{defn}

\begin{rem}[the fixed-point/curve model, cross-checked against an independent count]
$s_{ij}$ is the reflection in the root $e_i-e_j$, acting on the $T$-fixed points
of $LG(N,2N)$; the associated $T$-curve has weight $w_i/w_j$, so ``$\equiv$ along
that curve'' means congruence modulo $1-w_i/w_j$, a unit multiple of $t$.  The
same model assigns to $e_a+e_b$ the map $\bM\mapsto\bM\sm\{a,b\}$ (when
$\{a,b\}\subseteq\bM$) and to $2e_m$ the map $\bM\mapsto\bM\sm\{m\}$ --- which is
precisely the shape of (QC) and (SRC) in \texttt{2026-10-17-c2}, so the model is
consistent with what is already proved.  Independent check of the \emph{count}:
at $N=2$ the model predicts $1+1+2+2=6$ curves over the four reflections of
$W(C_2)$, and Rick's $LG(2,4)$ GKM checker (UID 824, a different author's
instrument) tests exactly six congruences.
\end{rem}

%=====================================================================
\section{Lemma A1: the $A$-type multiplier is $1$, by symmetry}
\label{sec:A1}

\begin{lem}[$u$-symmetry]\label{lem:sym}
For all $N,n,r$ the normalised weight function
$\widetilde{W}^X_{N,n,r}(u\mid w\mid\bI)$ is a symmetric function of
$u_1,\dots,u_n$, for $X=B$ and $X=C$.
\end{lem}

\begin{proof}
Motegi's Motherfunction (eq.\ \texttt{Motherfunction}, L474--487) is
\[
   F_{N,n,r}(u\mid w\mid\bI)
   =\frac{1}{(n-r)!}\Bigl(\prod_{k=1}^{n}u_k^{-(N+n)}\Bigr)
     \sum_{\sigma\in S_n}T\bigl(u_{\sigma(1)},\dots,u_{\sigma(n)}\bigr),
\]
for a fixed expression $T$ depending on $w,\bI,\beta,\gamma$.  A sum of
$T\circ\sigma$ over the \emph{full} symmetric group is symmetric whatever $T$ is,
and the prefactor $\prod_k u_k^{-(N+n)}$ is symmetric; hence $F$ is symmetric
in~$u$.  The type-$C$ normaliser $\mathcal{E}^C$ (eq.\
\texttt{typeCnormalizationfactor}, L1057--1062) is built from
$\prod_{j,k}(1-qu_k/u_j)$, $\prod_{j,k}(1-qu_jw_k)$, $\prod_{j<k}(1-qu_ju_k)$,
a $u$-free factor and $(1-q)^{n-r}$ --- each symmetric in $u$.  For type $B$,
$W^B$ carries the extra symmetric prefactor $\prod_k(1+1/u_k)^{-1}$ and
$\beta,\gamma=1,-1$, and $\mathcal{E}^B$ (L1361--1367) adds
$\prod_j(1-qu_j)$ and $\prod_j(1-qu_j^2)$, again symmetric.  Therefore
$\widetilde{W}^X=W^X/\mathcal{E}^X$ is symmetric in $u$.
\end{proof}

\begin{lem}[regularity on $A$-type divisors]\label{lem:reg}
$\At_{\bM,\bI}\in\mathcal{O}_{ij}$ for every $\bM,\bI\in P$ and every $i<j$; the
same holds for type $B$.
\end{lem}

\begin{proof}
Two sources of a possible pole must be excluded.

\emph{(a) The normaliser.}  $\At=W^X\cdot\mathcal{E}^X{}^{-1}$, so poles can come
only from zeros of the numerator of $\mathcal{E}^X$, namely from
$\prod_{j,k}(1-qu_k/u_j)$, $\prod_{j,k}(1-qu_jw_k)$ and (type $B$)
$\prod_j(1-qu_j)$.  At $u=w_\bM$ these vanish on $\{w_b=qw_a\}$,
$\{qw_aw_b=1\}$ and $\{qw_a=1\}$ respectively.  Each carries a $q$; none of them
contains the generic point of $\{w_i=w_j\}$, since $q$ is transcendental over
$\mathbb{C}(w)$.  So $\ord(\mathcal{E}^X)=0$ and the normaliser contributes
nothing.

\emph{(b) The Motherfunction.}  In $T(v_1,\dots,v_n)$ the only factors with a
denominator are
\[
   \frac{qv_a-v_b}{v_a-v_b}\quad (a<b\le r)
   \qquad\text{and}\qquad
   \frac{qv_b-v_a}{v_b-v_a}\quad (a\le r<b),
\]
one for each slot pair $\{a,b\}$ with $\min(a,b)\le r$; each is a simple pole on
$\{v_a=v_b\}$.  Fix two distinct elements $p\ne p'$ of $\bM$, so $u_p=w_p$ and
$u_{p'}=w_{p'}$ are distinct variables.  For $\sigma\in S_n$ let
$a:=\sigma^{-1}(p)$, $b:=\sigma^{-1}(p')$ be the slots carrying them, and put
\[
   \sigma':=\sigma\circ(a\,b).
\]
Since $\sigma$ is a bijection, $a\ne b$ are well defined, so
$\sigma\mapsto\sigma'$ is a fixed-point-free involution of $S_n$.  I claim
\[
   \Res_{w_p=w_{p'}}\Bigl[\,T(u_\sigma)+T(u_{\sigma'})\,\Bigr]=0 .
\]
Write $x:=u_{\sigma(a)}$, $y:=u_{\sigma(b)}$, so $\{x,y\}=\{w_p,w_{p'}\}$.  The
two terms differ \emph{only} by exchanging the values in slots $a$ and $b$.
Every factor other than the one attached to the pair $\{a,b\}$ is regular at
$x=y$: the pair $\{a,b\}$ is the unique slot pair whose two values are
$w_p,w_{p'}$, because $\sigma$ is injective.  Moreover at $x=y$ slots $a$ and $b$
carry equal values, so each of those other factors takes \emph{literally the same
value} in $T(u_\sigma)$ and in $T(u_{\sigma'})$; call that common value $C$.
For the dangerous factor, if $a<b\le r$ then
\[
  \Res_{x=y}\frac{qx-y}{x-y}=x(q-1),
  \qquad
  \Res_{x=y}\frac{qy-x}{y-x}=-\,y(q-1)=-\,x(q-1),
\]
and if $a\le r<b$ the same two values occur with the roles exchanged.  Either way
the two residues are $C\cdot x(q-1)$ and $-C\cdot x(q-1)$, and they cancel.
Summing over the $n!/2$ involution pairs, $\sum_\sigma T(u_\sigma)$ has no pole
on $\{w_p=w_{p'}\}$.  If $\min(a,b)>r$ the pair carries no such factor and both
terms are already regular.

Taking $\{p,p'\}=\{i,j\}$ when $\{i,j\}\subseteq\bM$ settles the only case in
which $w_i-w_j$ can occur as a difference of two $u$-arguments at all.  If
$|\bM\cap\{i,j\}|\le1$ then $w_i-w_j$ is never of the form $u_a-u_b$, and (b) is
vacuous.  Together with (a), $\ord(\At_{\bM,\bI})\ge0$.
\end{proof}

\begin{lem}[A1: the $A$-type GKM congruence, with multiplier exactly $1$]\label{lem:A1}
For all $N$, all $i<j$, all $\bM\in P$ and all $\bI\in P$,
\[
   \At_{\bM,\bI}\;\equiv\;\At_{s_{ij}\bM,\,\bI}\pmod{w_i-w_j},
\]
for type $C$ and for type $B$.  Equivalently $\rho^A=1$: no renormalisation is
needed, and none can change this.
\end{lem}

\begin{proof}
If $|\bM\cap\{i,j\}|\neq1$ then $s_{ij}\bM=\bM$ and there is nothing to prove.
So assume without loss of generality $i\in\bM$, $j\notin\bM$, and put
$\bM^{\circ}:=\bM\sm\{i\}$, so $s_{ij}\bM=\bM^{\circ}\sqcup\{j\}$.  In
particular $|s_{ij}\bM|=|\bM|$: the reflection \textbf{preserves the cardinality}
and therefore the integers $n=|\bM|$ and $r=|\bI|$ are the same on both sides.

By Lemma~\ref{lem:reg} both sides lie in $\mathcal{O}_{ij}$, so both may be
reduced modulo $t=w_i-w_j$, i.e.\ evaluated under the substitution
$\varepsilon:w_j\mapsto w_i$.  Write
$G(u_1,\dots,u_n):=\widetilde{W}^X_{N,n,r}(u\mid w\mid\bI)$, a function of $n$
arguments with the alphabet $w$ and the label $\bI$ as parameters.  Then
\[
   \At_{\bM,\bI}=G\bigl(w_i,(w_m)_{m\in\bM^{\circ}}\bigr),
   \qquad
   \At_{s_{ij}\bM,\bI}=G\bigl(w_j,(w_m)_{m\in\bM^{\circ}}\bigr),
\]
where by Lemma~\ref{lem:sym} the order of the arguments is immaterial.  Apply
$\varepsilon$ to both.  It acts on the parameters $w$ identically in the two
cases, and on the argument tuples it gives
\[
   \varepsilon\bigl(w_i,(w_m)_{m\in\bM^{\circ}}\bigr)
   =\bigl(w_i,(\varepsilon w_m)_{m\in\bM^{\circ}}\bigr)
   =\varepsilon\bigl(w_j,(w_m)_{m\in\bM^{\circ}}\bigr),
\]
the middle equality because $j\notin\bM^{\circ}$, so $\varepsilon$ does not touch
the tail, and $\varepsilon(w_j)=w_i$.  The two argument multisets therefore
become \emph{equal}, and the two sides become the same expression.
\end{proof}

\begin{rem}[why the pair divisors behaved differently --- the whole mechanism]
\label{rem:mech}
Lemma~\ref{lem:A1} is short, and the reason it is short is the content.  Compare
the three root families by what they do to $\bM$:
\[
\begin{array}{lll}
\text{root} & \text{action on }\bM & \text{effect on }|\bM| \\\hline
e_i-e_j & \bM\mapsto\bM\triangle\{i,j\} & \textbf{preserved} \\
e_a+e_b & \bM\mapsto\bM\sm\{a,b\} & \text{drops by }2 \\
h\cdot e_m & \bM\mapsto\bM\sm\{m\} & \text{drops by }1
\end{array}
\]
Only in the $A$-type case is the collision a \emph{permutation of the argument
multiset} $w_\bM$; $u$-symmetry then supplies the congruence with multiplier $1$,
for free and for every $h$.  In the other two cases the number of arguments
changes, symmetry says nothing, and the multiplier has to be \emph{computed} from
a genuine stability/collapse relation.

\textbf{What this does and does not explain.}  It explains why the $A$-type case
cost nothing: it is the one family whose multiplier is forced without computing
it.  It does \emph{not} predict which families carry poles, and I want to be
precise about that, because the stronger reading is tempting and false.  Changing
the arity is \emph{necessary but not sufficient} for a nontrivial multiplier: the
boundary root $h\cdot e_m$ also drops $|\bM|$, and its multiplier, once computed,
is \textbf{exactly $1$} (Theorem (SRC) of \texttt{2026-10-17-c2}).  So the
boundary family is the counterexample to ``arity changes $\Rightarrow$
$\rho\ne1$''.  The honest statement is: arity preserved $\Rightarrow$ $\rho=1$
by symmetry; arity changed $\Rightarrow$ no information, go and compute, and the
two computations landed differently ($\rho^X\ne1$ for $e_a+e_b$, $\rho=1$ for
$h\cdot e_m$).

Given the multipliers, though, the pole census follows exactly.  The product of
two basis elements collides with multiplier $\rho^2$ while their span collides
with $\rho$, so the obstruction is $\rho^2-\rho=\rho(\rho-1)$, which vanishes
precisely when $\rho\in\{0,1\}$.  On the pair divisors $\rho=g\ne1$ and there are
poles ($13/115$ at $N=3$); on the boundary divisors $\rho=1$ and there are none;
on the $A$-type divisors $\rho^A=1$ and $1^2-1=0$.  \emph{The poles were never
about the $A$-direction, and the renormalisation $\psi$ was never needed there.}
\end{rem}

%=====================================================================
\section{Theorem A: the span is exactly the GKM lattice}
\label{sec:thmA}

Fix $i<j$ and abbreviate $\mathcal{O}:=\mathcal{O}_{ij}$, $s:=s_{ij}$,
$t:=w_i-w_j$, $k:=k_{ij}$.

\begin{defn}
\[
  \Gkm:=\bigl\{f\in\mathcal{O}^{P}\;:\;f(\bM)\equiv f(s\bM)\ (\mathrm{mod}\ t)
         \ \ \forall\,\bM\in P\bigr\},
  \qquad
  \Lat:=\bigoplus_{\bI\in P}\mathcal{O}\,\phi_\bI .
\]
\end{defn}

\begin{lem}[$\Gkm$ is a ring]\label{lem:ring}
$\Gkm$ is an $\mathcal{O}$-subalgebra of $\mathcal{O}^{P}$ containing $1$.
\end{lem}

\begin{proof}
The constant function $1$ satisfies the congruences.  $\Gkm$ is an
$\mathcal{O}$-submodule because the conditions are $\mathcal{O}$-linear.  For
closure under multiplication, for $f,g\in\Gkm$ and any $\bM$,
\[
  (fg)(\bM)-(fg)(s\bM)
  = f(\bM)\bigl[g(\bM)-g(s\bM)\bigr]+\bigl[f(\bM)-f(s\bM)\bigr]g(s\bM)
  \in t\mathcal{O},
\]
since all four of $f(\bM),g(s\bM)$ and the two brackets lie in $\mathcal{O}$ and
the brackets lie in $t\mathcal{O}$.
\end{proof}

\begin{lem}[colength of $\Gkm$]\label{lem:colenG}
$\operatorname{length}_{\mathcal{O}}\bigl(\mathcal{O}^{P}/\Gkm\bigr)=2^{N-2}$.
\end{lem}

\begin{proof}
Choose one representative $\bM_\alpha$ from each of the $2^{N-2}$ two-element
$s$-orbits $\alpha$ and consider
\[
   \delta:\mathcal{O}^{P}\longrightarrow k^{\,2^{N-2}},
   \qquad
   \delta(f)_\alpha:=\bigl[f(\bM_\alpha)-f(s\bM_\alpha)\bigr]\bmod t .
\]
This is $\mathcal{O}$-linear, and $\ker\delta=\Gkm$: the congruence at
$s\bM_\alpha$ is the negative of the one at $\bM_\alpha$, and at an $s$-fixed
point it is $0\equiv0$.  It is surjective: for a target vector
$(\bar c_\alpha)$, lift each $\bar c_\alpha$ to $c_\alpha\in\mathcal{O}$ and take
$f=\sum_\alpha c_\alpha\mathbf{1}_{\bM_\alpha}$, whose $\delta$ is
$(\bar c_\alpha)$ because the points $\bM_\alpha,s\bM_\alpha$ are pairwise
distinct across orbits.  Hence
$\mathcal{O}^{P}/\Gkm\cong k^{2^{N-2}}$, of length $2^{N-2}$.
\end{proof}

\begin{lem}[colength of $\Lat$]\label{lem:colenL}
$\operatorname{length}_{\mathcal{O}}\bigl(\mathcal{O}^{P}/\Lat\bigr)
=\sum_{\bI\in P}\ord\At_{\bI,\bI}=2^{N-2}$.
\end{lem}

\begin{proof}
By Lemma~\ref{lem:reg} the matrix $\Phi=(\At_{\bM,\bI})$ has entries in
$\mathcal{O}$, and $\Lat$ is exactly its column span, i.e.\ the image of
$\Phi:\mathcal{O}^{P}\to\mathcal{O}^{P}$.  Since $\det\Phi\ne0$, the map is
injective, and over a discrete valuation ring the colength of the image of an
injective endomorphism of a free module of finite rank is the valuation of its
determinant (Smith normal form).  By (M3)+(M4),
$\det\Phi=\pm\prod_\bI\At_{\bI,\bI}$, so
$\operatorname{length}(\mathcal{O}^{P}/\Lat)=\sum_\bI\ord\At_{\bI,\bI}$.

Now evaluate $\ord\At_{\bI,\bI}$ from \eqref{eq:diag}, factor by factor, by
asking which factors contain the generic point of $\{w_i=w_j\}$.  The factors
$\bigl(1-qw_cw_d\bigr)$ and $\bigl(1-qw_c/w_{i_l}\bigr)$ carry a $q$ and so do
not (same argument as Lemma~\ref{lem:reg}(a)).  The factor
$\bigl(1-1/(w_{i_a}w_{i_b})\bigr)$ cuts the pair divisor $\{w_{i_a}w_{i_b}=1\}$
(or, at $a=b$, the boundary divisor $\{w_{i_a}^2=1\}$), an irreducible divisor
distinct from $\{w_i=w_j\}$, so it does not either.  The only remaining factors
are the $A$-type ones
\[
  \Bigl(1-\frac{w_c}{w_{i_l}}\Bigr),\qquad c\in\bI^{c},\ i_l\in\bI,\ i_l>c,
\]
and such a factor vanishes on $\{w_i=w_j\}$ if and only if $\{c,i_l\}=\{i,j\}$.
Given $i<j$ and the constraint $i_l>c$, this forces $c=i$ and $i_l=j$.  Each such
factor is simple and occurs at most once.  Therefore
\[
  \ord\At_{\bI,\bI}=
  \begin{cases}1,& i\notin\bI\ \text{and}\ j\in\bI,\\ 0,&\text{otherwise,}\end{cases}
  \qquad\text{so}\qquad
  \sum_{\bI\in P}\ord\At_{\bI,\bI}=\#\{\bI: i\notin\bI,\ j\in\bI\}=2^{N-2}. \qedhere
\]
\end{proof}

\begin{lem}[the $A$-type blocks of the diagonal do not know the type]\label{lem:typeB}
Lemma~\ref{lem:colenL} holds verbatim for type $B$: for all $\bI$,
$\ord_{(w_i-w_j)}\At^{B}_{\bI,\bI}=[\,i\notin\bI\ \text{and}\ j\in\bI\,]$ and
$\sum_\bI\ord\At^{B}_{\bI,\bI}=2^{N-2}$.
\end{lem}

\begin{proof}
Motegi gives the same root-product shape in both types ---
$\prod_{\mathrm{Inv}(\lambda)}(1-e^{-\alpha})\prod_{T(\lambda)}(1-qe^{+\alpha})$,
eq.\ \texttt{normalizationWC} (L1078) and eq.\ \texttt{normalizationWB} (L1386)
--- and writes out $\mathrm{Inv}$ and $T$ explicitly in each case: type $C$ at
L997--1010 and L1017--1030, type $B$ at L1307--1320 and L1322--1340.  Read
first-hand this session, the two decompositions are
\begin{align*}
\mathrm{Inv}(\lambda)\ \text{type }C &=
   \{2t_k\}_{k\in\bI}\ \cup\ \{t_j{+}t_k\}_{j<k;\,j,k\in\bI}\ \cup\
   \mathbf{\{t_k{-}t_j\}_{j<k;\ j\notin\bI,\,k\in\bI}}, \\
\mathrm{Inv}(\lambda)\ \text{type }B &=
   \phantom{2}\{t_k\}_{k\in\bI}\ \cup\ \{t_j{+}t_k\}_{j<k;\,j,k\in\bI}\ \cup\
   \mathbf{\{t_k{-}t_j\}_{j<k;\ j\notin\bI,\,k\in\bI}}, \\[1mm]
T(\lambda)\ \text{type }C &=
   \{2t_k\}_{k\notin\bI}\ \cup\ \{t_j{+}t_k\}_{j<k;\,j,k\notin\bI}\ \cup\
   \mathbf{\{t_k{-}t_j\}_{j<k;\ j\in\bI,\,k\notin\bI}}, \\
T(\lambda)\ \text{type }B &=
   \phantom{2}\{t_k\}_{k\notin\bI}\ \cup\ \{t_j{+}t_k\}_{j<k;\,j,k\notin\bI}\ \cup\
   \mathbf{\{t_k{-}t_j\}_{j<k;\ j\in\bI,\,k\notin\bI}}.
\end{align*}
The $A$-type members (bold) are \emph{the same subsets of the same root lattice},
and the corresponding product blocks are correspondingly identical in the two
papers' own notation:
\[
  \prod_{j\in\bI^{c}}\prod_{l:\,i_l>j}\bigl(1-w_{i_l}^{-1}w_j\bigr)
  \quad\text{from }\mathrm{Inv},
  \qquad
  \prod_{j\in\bI^{c}}\prod_{l:\,i_l<j}\bigl(1-q\,w_{i_l}^{-1}w_j\bigr)
  \quad\text{from }T .
\]
The \emph{entire} type-dependence sits on the boundary root: $2t_k$ versus
$t_k$, i.e.\ $(1-w_{i_j}^{-2})$ versus $(1-w_{i_j}^{-1})$ in $\mathrm{Inv}$ and
$(1-qw_k^2)$ versus $(1-qw_k)$ in $T$ --- which is also why type $C$ can absorb
its $j=k$ term into $\prod_{j\le k}$ while type $B$ carries a separate
$\prod_j(1-w_{i_j}^{-1})$ and a strict $\prod_{j<k}$.  Neither
$\{w_{i_j}=1\}$ nor $\{qw_k=1\}$ contains the generic point of $\{w_i=w_j\}$, so
the factor-by-factor count in the proof of Lemma~\ref{lem:colenL} is unchanged,
and only the bold blocks contribute.
\end{proof}

\begin{rem}[the type enters only through $h$, for the third time]
Lemma~\ref{lem:typeB} is the same phenomenon as the previous paper's thesis and
as this one's Remark~\ref{rem:mech}: type $B$ and type $C$ differ \emph{only} in
the boundary root $h\cdot e_m$, and every statement about the $A$-type direction
is therefore type-free.  It is worth recording that I did \emph{not} obtain this
by assuming it.  I was about to file the type-$B$ arm as a gap resting on a
measured colength, and went to check a \emph{different} locator --- the one for
Motegi's order --- which turned out to be wrong; finding the right one put me in
the neighbourhood of \texttt{normalizationWB}, which I then read.  Reusing
\eqref{eq:diag} for type $B$ \emph{without} reading L1307--1340 would have been
the exact shape of the $\mathcal{E}^B$ error of 2026-10-17 (a type-$C$ formula
applied to type $B$, green and wrong); the difference between that error and this
lemma is four lines of reading, not an inference.
\end{rem}

\begin{thm}[the span \emph{is} the $A$-type GKM lattice]\label{thm:A}
For every $N\ge2$, every $i<j$, and \emph{both} types $B$ and $C$,
\[
   \Lat_{ij}\;=\;\Gkm_{ij}.
\]
\end{thm}

\begin{proof}
$\Lat\subseteq\Gkm$: each $\phi_\bI$ lies in $\mathcal{O}^P$ by
Lemma~\ref{lem:reg} and satisfies the congruences by Lemma~\ref{lem:A1}, and
$\Gkm$ is an $\mathcal{O}$-module.  Both have colength $2^{N-2}$ in
$\mathcal{O}^P$ by Lemmas~\ref{lem:colenG} and~\ref{lem:colenL} (type $C$), resp.\
Lemma~\ref{lem:typeB} (type $B$), so
\[
  \operatorname{length}(\Gkm/\Lat)
  =\operatorname{length}(\mathcal{O}^P/\Lat)-\operatorname{length}(\mathcal{O}^P/\Gkm)
  =2^{N-2}-2^{N-2}=0,
\]
the subtraction being legitimate because both lengths are finite.  A module of
length $0$ is zero, so $\Lat=\Gkm$.
\end{proof}

\begin{rem}[the two $2^{N-2}$'s are not the same count wearing one name]
On the GKM side $2^{N-2}$ counts \emph{two-element orbits of $s_{ij}$ on $P$}.
On the span side it counts \emph{labels $\bI$ with $i\notin\bI$, $j\in\bI$}.
These index sets are genuinely different --- an orbit is an unordered pair of
points, a label is a single subset --- and the bijection between them
($\{\bM,s\bM\}\mapsto$ the member containing $j$, intersected with the
complement of $\{i\}$) is not used anywhere in the proof.  The two numbers are
computed by different routes and then compared, which is the only way the
comparison carries information.
\end{rem}

%=====================================================================
\section{Theorem B: no condition on $\chi$}
\label{sec:thmB}

\begin{defn}
For $\chi\in F(v)$ put $\Psi_\chi(\bM):=\prod_{m\in\bM}\chi(w_m)$ and
$\Ah^{\chi}_{\bM,\bI}:=\At_{\bM,\bI}\,\Psi_\chi(\bM)$, i.e.\
$\hat\phi^{\chi}_\bI=\Psi_\chi\cdot\phi_\bI$ as functions on $P$.  Call $\chi$
\textbf{$ij$-admissible} if $\chi(w_m)\in\mathcal{O}_{ij}^{\times}$ for every
$m\in[N]$, i.e.\ $\chi(w_m)$ is regular and non-vanishing at the generic point of
$\{w_i=w_j\}$.
\end{defn}

\begin{lem}[$\Psi_\chi$ is a unit of $\Gkm$ --- and this is why it carries no information]
\label{lem:unit}
If $\chi$ is $ij$-admissible then $\Psi_\chi\in\Gkm^{\times}$, i.e.\ both
$\Psi_\chi$ and $\Psi_\chi^{-1}$ lie in $\Gkm$.
\end{lem}

\begin{proof}
$\Psi_\chi(\bM)$ and $\Psi_\chi(\bM)^{-1}$ are finite products of the units
$\chi(w_m)^{\pm1}$, hence lie in $\mathcal{O}^{\times}\subset\mathcal{O}$.  For
the congruence take $\bM$ with $i\in\bM$, $j\notin\bM$ and
$\bM^{\circ}=\bM\sm\{i\}$; then
\[
  \Psi_\chi(\bM)-\Psi_\chi(s\bM)
  =\Psi_\chi(\bM^{\circ})\bigl[\chi(w_i)-\chi(w_j)\bigr]\in t\mathcal{O},
\]
because $\chi(w_i)-\chi(w_j)\in t\mathcal{O}$: writing
$\chi=\chi_1/\chi_2$ with $\chi_2(w_i),\chi_2(w_j)$ units, the difference is
$\bigl(\chi_1(w_i)\chi_2(w_j)-\chi_1(w_j)\chi_2(w_i)\bigr)/\chi_2(w_i)\chi_2(w_j)$
and the numerator is a polynomial expression vanishing under $w_j\mapsto w_i$.
The identical computation with $\chi^{-1}$ in place of $\chi$ gives
$\Psi_\chi^{-1}\in\Gkm$.
\end{proof}

\begin{rem}[Task 1 of the brief, and its exact epistemic status]\label{rem:task1}
Lemma~\ref{lem:unit} is the observation the brief flagged at Task~1 and
instructed me to grade honestly, so I do.  \textbf{It carries zero bits about
$\psi$.}  The proof uses nothing about $\chi$ beyond admissibility; it goes
through verbatim for $\chi(v)=v$, for $\chi(v)=1$, for
$\chi(v)=\prod(\text{anything regular and nonvanishing})$.  The pre-registered
``non-invariant'' control $\prod_{m\in\bM}w_m+|\bM|$ failed to fire at the
2026-10-18 wake for exactly this reason, and that silent control was the finding:
the invariance is a property of the \emph{shape} $\prod_{m\in\bM}f(w_m)$, not of
$\psi$.  Measured there: $0$ violations in $96$ instances at $N=4$, $h=2$, with
the off-divisor arm firing $48/48$.

And yet it is logically indispensable here, which is the part worth recording.
Having zero discriminating power between $\Ah$ and $\At$ is precisely what makes
Theorem~\ref{thm:B} hold for the \emph{whole} family rather than for $\psi$
alone.  A fact can be useless as evidence and load-bearing as a hypothesis; the
error would be to let the second tempt me into reporting the first as support
for $\psi$.  It is not.  Cf.\ the one-line version in \S\ref{sec:branch}.
\end{rem}

\begin{thm}[every structure constant is regular on every $A$-type divisor]\label{thm:B}
Fix $i<j$ and let $\Psi\in\Gkm_{ij}^{\times}$ be \emph{any} invertible element of
the ring $\Gkm_{ij}$ --- not necessarily a point-indexed product.  Let
$\hat\phi_\bI:=\Psi\,\phi_\bI$ and let $\hat c^{\nu}_{\lambda\mu}$ be the
resulting structure constants.  Then for all $\lambda,\mu,\nu$,
\[
   \hat c^{\nu}_{\lambda\mu}\in\mathcal{O}_{ij}.
\]
\end{thm}

\begin{proof}
Let $\hat\Lat:=\bigoplus_\bI\mathcal{O}\hat\phi_\bI=\Psi\cdot\Lat$.  By
Theorem~\ref{thm:A} and because $\Psi$ is a unit of the ring $\Gkm$
(Lemma~\ref{lem:ring}),
\[
  \hat\Lat=\Psi\,\Lat=\Psi\,\Gkm=\Gkm ,
\]
so $\hat\Lat$ is a ring.  Hence $\hat\phi_\lambda\hat\phi_\mu\in\hat\Lat$, and
expanding it in the $\mathcal{O}$-basis $\{\hat\phi_\nu\}$ of $\hat\Lat$ gives
coefficients in $\mathcal{O}$.  These coefficients are the
$\hat c^{\nu}_{\lambda\mu}$, by uniqueness of the expansion in the $F$-basis
$\{\hat\phi_\nu\}$ of $F^P$.
\end{proof}

\begin{cor}[the point-product family]\label{cor:family}
If $\chi$ is $ij$-admissible for all $i<j$, then every structure constant of
$\Psi_\chi\At$ is regular on every $A$-type divisor.  In particular this holds
for $\At$ itself ($\chi=1$) and for $\Ah=\Psi_\psi\At$, and the $A$-type
behaviour is the same for every member of the family.
\end{cor}

\begin{proof}
Lemma~\ref{lem:unit} puts $\Psi_\chi$ in $\Gkm_{ij}^{\times}$ for each $i<j$;
apply Theorem~\ref{thm:B} for each.
\end{proof}

\begin{rem}[the hypothesis is $\Gkm^{\times}$-membership, and \emph{not} either of
the two things it is easy to mistake it for]\label{rem:hyp}
I stated Theorem~\ref{thm:B} for the point-product family first and then
lesioned it (V7), which moved the hypothesis twice.
\begin{itemize}
\item It is \textbf{not} ``$\Psi$ is a point-indexed product''.  The lesion
  $\Psi=1+(w_1-w_2)$ on the single point $\bM=\{1\}$ and $1$ elsewhere is
  \emph{not} of the form $\prod_{m\in\bM}\chi(w_m)$, and the conclusion holds
  anyway (A-type set empty) --- because $1+(w_1-w_2)\equiv1$, so that $\Psi$ lies
  in $\Gkm^{\times}$ regardless.  A surviving lesion is a stronger theorem for
  free, and this is why Theorem~\ref{thm:B} is now stated for $\Gkm^{\times}$
  with the product family demoted to Corollary~\ref{cor:family}.
\item It is \textbf{not} ``$\Psi$ is a pointwise unit'' either.  The lesion
  $\Psi=2$ on $\bM=\{1\}$ and $1$ elsewhere is a unit at every point, with a
  unit inverse, and it produces \textbf{six} $A$-type poles --- because
  $2\not\equiv1$, so it fails the congruence and is not in $\Gkm$ at all.
\end{itemize}
The two lesions differ only in whether the perturbation is $\equiv1$ modulo the
uniformiser, and they land on opposite sides.  That pair is the sharp control for
this hypothesis: together they show $\Gkm^{\times}$ is neither too strong (the
first survives) nor too weak (the second breaks).
\end{rem}

\begin{cor}[the renormalisation-invariance statement]\label{cor:inv}
The lattice $\Lat_{ij}$ is independent of $\chi$ within the admissible family:
$\bigoplus_\bI\mathcal{O}_{ij}\,\Psi_\chi\phi_\bI=\Gkm_{ij}$ for every
$ij$-admissible $\chi$.  Consequently the $A$-type divisors impose
\textbf{no condition} on $\chi$.
\end{cor}

\begin{rem}[the exact complement of Theorem C of \texttt{2026-10-17-c2}]
Theorem C of the previous paper determined the admissible renormalisations by
\emph{two} conditions and showed both are necessary: (COB)
$\chi(x)\chi(1/x)=\rho^{-1}|_H$ for the pair divisors, and (SR)
$\chi(\varepsilon)=1$ for $\varepsilon^h=1$ for the boundary divisors.  It did
not know whether a \emph{third} condition was waiting at the $A$-type divisors.
Corollary~\ref{cor:inv} says there is none: (COB) and (SR) are the complete list.
This sharpens Theorem C from ``these two conditions suffice for the two families
examined'' to ``these two conditions are the whole story across all three
families of positive roots'', and it is the statement Q510 needs --- a
characterisation of $\Ah$ cannot be read off the $A$-direction, because the
$A$-direction is blind to the renormalisation.  Note the scope limit inherited
from the previous paper's Theorem C: the family considered is
\emph{point-indexed and diagonal}, $\Psi_\chi(\bM)=\prod_{m\in\bM}\chi(w_m)$.
Koncki's discrepancy at $N=2$ is unipotent \emph{off-diagonal in label space}
(Rick, UID 824), so it lies outside this family and outside
Corollary~\ref{cor:inv}'s reach.
\end{rem}

%=====================================================================
\section{The pre-registered branch, and what it cost me}
\label{sec:branch}

The brief wrote the branch before the measurement, and I quote it so the grading
is checkable:

\begin{quote}\small
``\textbf{PRE-REGISTERED BRANCH, written before measuring.} If $\At$ has
\textbf{no} $A$-type poles, then Task 2 is vacuous too and \textbf{gap 2 was
never a gap}.  Report that plainly and stop; do not dress a vacuity as a result.
I expect $\At$ \emph{does} have them (the bulk is a type-$A$-symmetric model and
$\mathcal{E}^C$ is a boundary normaliser, so nothing cancels the $A$-type
collisions), but I have been wrong about exactly this shape of expectation twice
this week, both times on a normaliser.''
\end{quote}

\textbf{The branch fired on the arm I did not expect, for the third time this
week.}  $\At$ has no $A$-type poles --- not in the basis (all $35$ nonzero
entries at $N\le3$ are Laurent polynomials, \S\ref{sec:verif}), and not in the
structure constants (the $13/115$ non-Laurent baseline constants have
\emph{pair-divisor denominators only}; the $A$-type set is empty).  So there was
no $A$-type defect for $\Psi$ to create or destroy, and \textbf{gap 2 read as a
difference question --- ``does $\Psi$ break the $A$-direction?'' --- is
vacuous.}  I say that plainly, as instructed.

Where my stated reason went wrong is worth naming precisely, because the error
was not in the conclusion but in the \emph{direction of the inference}.  I wrote
``the bulk is a type-$A$-symmetric model \dots\ so nothing cancels the $A$-type
collisions.''  The symmetry is real and I had the right object; I then used it to
predict that the collisions \emph{survive}.  It is the very thing that makes them
\emph{harmless} --- $u$-symmetry is exactly the hypothesis of
Lemma~\ref{lem:A1}.  I had the mechanism and read its sign backwards.

The two readings of gap 2 therefore come apart, and only one of them is empty:
\begin{itemize}
\item \emph{As a difference question} (``is $\Ah$ worse than $\At$ here?''):
  vacuous, and provably so --- Corollary~\ref{cor:inv} says no admissible
  renormalisation changes the $A$-type lattice at all.
\item \emph{As a regularity question} (``are the structure constants regular on
  $\{w_i=w_j\}$?''): \textbf{not vacuous}.  It was unproved before today for
  either basis, at any $N$; \texttt{2026-10-17-c2} proved regularity only on the
  pair and boundary divisors, and its Theorem~B's unit argument gives nothing
  here (a unit cannot cancel a pole in the basis collapse, but the structure
  constants \emph{reshuffle}).  Theorem~\ref{thm:B} settles it for all $N$ and
  for the whole family, and Theorem~\ref{thm:A} says something strictly stronger
  than regularity: the span is not merely \emph{contained in} a ring, it
  \emph{equals} the GKM lattice.
\end{itemize}

The honest summary: the session's measurement was vacuous and its theorem was
not.  A vacuity has a mechanism, and in this case the mechanism
(Remark~\ref{rem:mech}) also explains the \emph{non}-vacuous half of the previous
paper --- why the pair divisors, and only they, carried poles.

%=====================================================================
\section{Verification}
\label{sec:verif}

\begin{sloppypar}
All computations are pure Python + SymPy (SageMath is not installed in this
container) in \texttt{proofs/code-1018-atype/}.  The engine
\texttt{q484\_engine.py} exists in four places in the tree
(\texttt{code-1018-atype/}, \texttt{code-1017c2-coboundary/},
\texttt{code-1017-nonvanH/}, \texttt{reviews/scratch-20261017c2/}).  Checked as a
\emph{set} at the start of this session: all four byte-identical, md5
\texttt{8b37ab1b}, so the stale-copy warning carried in the brief is discharged.
Re-checked as a set after the locator fix described in (M3): all four
\texttt{f53c0258}.  I verify the \emph{set}, not one pair, because a drift check
run on one pair reports health for the copy it did not open.
\end{sloppypar}

\paragraph{V0. Root lengths, from first principles.}
Enumerating $B_n$ and $C_n$ for $n=2,3,4$ and classifying by $|\alpha|^2$: the
norms are $\{1,2\}$ in $B$ and $\{2,4\}$ in $C$, so $e_a+e_b$ is \emph{long} in
$B$ and \emph{short} in $C$, while $h\cdot e_m$ is \emph{short} in $B$ and
\emph{long} in $C$.  This is the measurement behind the nomenclature correction
to \texttt{2026-10-17-c2} (Remark 2 there, commit \texttt{7002d00}).

\paragraph{V1. The divisor classifier.}
\texttt{divisors.py} decides a factor's divisor class \emph{by vanishing} ---
substituting each candidate divisor's defining relation --- rather than by
matching the printed form, which would miss $1-w_2/w_1$ and $(w_1-w_2)\cdot(-1)$.
Positive controls: $11$ planted factors, each placed in the correct one of
\texttt{A}/\texttt{pair}/\texttt{bdy}/\texttt{qdef}/\texttt{other}/\texttt{unit},
$11/11$.  Crucially $(w_a-qw_b)$ and $(1-qw_aw_b)$ land in \texttt{qdef} and are
never merged into a geometric class: $\{w_a=qw_b\}$ and $\{w_a=w_b\}$ differ as
sets and coincide only at $q=1$.

\paragraph{V2. The basis is Laurent (\texttt{basis\_atype.py}).}
At $N=2$ ($10$ nonzero entries) and $N=3$ ($35$ nonzero entries), type $C$:
\emph{no non-monomial denominator occurs at all}, so in particular zero $A$-type
poles.  This is stronger than Lemma~\ref{lem:reg}, which claims only
$\ord\ge0$ at $A$-type divisors; I do not claim the stronger statement as proved.

\paragraph{V3. Lemma A1, with three controls (\texttt{atype\_gkm.py}).}
Exact rational $w$ with $w_j=w_i$ imposed \emph{exactly}, $q$ kept symbolic, so
each test is an identity of rational functions in $q$; two independent $w$-seeds.

\begin{center}\small
\begin{tabular}{llll}
\toprule
 & $N=2$ & $N=3$ & $N=4$ \\\midrule
(i,j,$\bM$,$\bI$) instances, both seeds & $8/8$ & $96/96$ & $768/768$ \\
of which trivial $0=0$ & $4$ & $54$ & $464$ \\
with a nonzero side & $4$ & $42$ & $304$ \\
$\rho^A=1$ and $\bI$-independent, every orbit & yes ($2$) & yes ($12$) & yes ($48$) \\\midrule
C1 off divisor ($w_j=\tfrac54w_i$) & fails $6$ & fails $62$ & fails $442$ \\
C2 planted $\rho=1+q$ & fails $4$ & fails $42$ & fails $304$ \\
C3 wrong partner $s_{ik}\bM$, $k\ne i,j$ & --- & fails $62$ & fails $456$ \\
\bottomrule
\end{tabular}
\end{center}

All three controls fire.  Each is recorded with the factor it is sensitive to:
C1 to the collision itself; C2 to the \emph{value} of the multiplier; C3 to
\emph{which point} the congruence pairs $\bM$ with.  C2's ``holds'' column is
$4$, $54$, $464$ --- \emph{exactly} the trivial-$0=0$ counts, i.e.\ the plant was
rejected on every one of the $4$, $42$, $304$ nonzero cells and structurally
spared only where $0=\rho\cdot0$ holds for any $\rho$.  Reporting the nonzero
column separately is what makes that green readable.

\paragraph{V4. The colength arithmetic (\texttt{diag\_order.py}).}
$\ord_{(w_i-w_j)}\At_{\bI,\bI}$ computed \emph{directly from the engine} --- the
$S_n$ symmetrisation --- and compared against the formula derived from the root
product \eqref{eq:diag}.  These are independent mechanisms; it was the same kind
of cross-check that caught the $\mathcal{E}^B$ error on 2026-10-17.
\begin{center}\small
\begin{tabular}{llll}
\toprule
 & $\sum_\bI\ord$ vs.\ $2^{N-2}$ & per-label $\ord=[i\notin\bI,j\in\bI]$ & (M4) engine $=$ root product \\\midrule
$N=2$, $C$, $1$ pair  & $1=1$ & all $4$ labels agree & $4/4$ \\
$N=3$, $C$, $3$ pairs & $2=2$ & all $8$ labels agree, every pair & $8/8$ \\
$N=4$, $C$, $6$ pairs & $4=4$ & all $16$ labels agree, every pair & $16/16$ \\
$N=2$, $B$, $1$ pair  & $1=1$ & all $4$ labels agree & n/a \\
$N=3$, $B$, $3$ pairs & $2=2$ & all $8$ labels agree, every pair & n/a \\
\bottomrule
\end{tabular}
\end{center}
The type-$B$ rows were run as a \emph{measurement}, because the engine has no
type-$B$ \texttt{diag\_rootproduct} and \eqref{eq:diag} must not be reused there
(that reuse is precisely the $\mathcal{E}^B$ error).  They have since become a
check on Lemma~\ref{lem:typeB}, which proves the type-$B$ count from Motegi's own
\texttt{normalizationWB}; see Remark~\ref{gap:B}.

\paragraph{V5. The residue mechanism of Lemma~\ref{lem:reg}
(\texttt{residue\_pairing.py}).}
For five strata $(\bM,\bI)$ with $\{1,2\}\subseteq\bM$: (T1) the full
$\sum_{\sigma\in S_n}$ is regular at $w_1=w_2$ in $5/5$; (T2) \emph{every
involution pair} $\{\sigma,\sigma'\}$ is \emph{already} regular, $5/5$ strata,
and the involution is fixed-point-free as claimed.  So the cancellation is
pairwise, which is the mechanism the proof asserts and not a global accident.

The control --- a partial sum over a $\sigma$-set \emph{not} closed under the
involution must have a pole --- fires in $2$ strata and is silent in $3$.  The
three silences are \textbf{empty domains, with a named cause in each case}, not a
blind instrument:
\begin{itemize}
\item $\bI=\emptyset$ gives $r=0$, and both ranges carrying a
  $\frac{qv-v'}{v-v'}$ factor are indexed by $k\le r$ --- so the stratum
  \emph{structurally deletes every factor the control tests}.  Measured: $0$ of
  $6$ summands have a pole.
\item $\bI=(2,3)$ and $\bI=(1,2)$: the factors
  $\Lambda_j\ni(v_j-w_1)(v_j-w_2)$ supply numerator zeros on $\{w_1=w_2\}$ that
  kill each summand's pole term-by-term.  Measured: $0$ of $6$, resp.\ $0$ of
  $2$, summands have a pole.
\end{itemize}
In the two strata where the tested factor \emph{is} present ($\bM=(1,2,3)$,
$\bI=(1,)$: $4$ of $6$ summands polar; $\bM=(1,2)$, $\bI=(1,)$: $2$ of $2$), the
control fires, $\ord=-1$ both times.

\paragraph{V6. Structure-constant census, denominators classified
(\texttt{sc\_census.py}).}
Same triangular solve as the previous paper's \texttt{arm.py}, so the $13/115$
and $0/115$ baselines are reproducible, but every non-monomial denominator factor
is classified.  Reported as \emph{sets} of triples $(\lambda,\mu,\nu)$, never as
counts.  Pre-registered in \texttt{PREDICTIONS.md} before running: \textbf{zero
$A$-type denominators in every arm}, including the four arms that produce
$57/115$ non-Laurent constants --- a single $(w_a-w_b)$ among those $57$ would
refute Theorem~\ref{thm:B}.
\begin{center}\small
\begin{tabular}{llrrcrrl}
\toprule
 & & \multicolumn{4}{c}{$N=2$ (complete)} & \multicolumn{2}{c}{$N=3$} \\
\cmidrule(lr){3-6}\cmidrule(lr){7-8}
arm & $\chi(v)$ & nz & non-Laur. & \textbf{\texttt{A}} & \texttt{pair}/\texttt{bdy} & non-Laur. & \textbf{\texttt{A}} \\\midrule
\texttt{id}    & $1$                   & $20$ & $1$  & $\mathbf{0}$ & $1/0$ & $13$ & $\mathbf{0}$ \\
\texttt{psi}   & $(1-qv^2)/(1-q)$      & $20$ & $0$  & $\mathbf{0}$ & $0/0$ & $0$  & $\mathbf{0}$ \\
\texttt{x2psi} & $v^2(1-qv^2)/(1-q)$   & $20$ & $0$  & $\mathbf{0}$ & $0/0$ & \multicolumn{2}{c}{not run} \\
\texttt{sq}    & $\bigl((1-qv^2)/(1-q)\bigr)^2$ & $20$ & $1$  & $\mathbf{0}$ & $1/0$ & \multicolumn{2}{c}{not run} \\
\texttt{h1}    & $(1-qv)/(1-q)$        & $20$ & $8$  & $\mathbf{0}$ & $1/8$ & \multicolumn{2}{c}{not run} \\
\texttt{unrm}  & $1-qv^2$              & $20$ & $8$  & $\mathbf{0}$ & $1/8$ & \multicolumn{2}{c}{not run} \\
\texttt{xpsi}  & $v(1-qv^2)/(1-q)$     & $20$ & $8$  & $\mathbf{0}$ & $0/8$ & \multicolumn{2}{c}{not run} \\
\texttt{x}     & $v$                   & $20$ & $8$  & $\mathbf{0}$ & $1/8$ & $\mathbf{57}$ & $\mathbf{0}$ \\\midrule
\texttt{atp}   & \emph{plant, see below} & $20$ & $14$ & $\mathbf{13}$ & $1/8$ & \multicolumn{2}{c}{\textsc{running}} \\
\bottomrule
\end{tabular}
\end{center}
At $N=2$ the sweep is complete over all eight admissible arms: $34$ non-Laurent
constants in total, and the \texttt{A} column is $0$ in every one of them --- $34$
separate opportunities to refute Theorem~\ref{thm:B}, all declined.

At $N=3$ the three arms that ran reproduce the previous paper's table exactly, and
the hardest one is decisive.  \texttt{id}: $115$ nonzero, $13$ non-Laurent, all
$13$ \texttt{pair}, \texttt{A}-set $\emptyset$.  \texttt{psi}: $0/115$.
\texttt{x}: $\mathbf{57}$ non-Laurent, of which $13$ \texttt{pair} and $57$
\texttt{bdy} --- and \texttt{A}-set $\emptyset$.  \texttt{x} is the arm that fails
\emph{both} (COB) and (SR), so it is where a third condition would have shown up if
there were one: \textbf{$57$ separate opportunities to refute
Theorem~\ref{thm:B} at $N=3$, all declined.}

\textbf{The cap.}  The five arms marked ``not run'' at $N=3$ were killed to hold
the machine at its four-core load, and \texttt{atp} at $N=3$ had not finished.  I
record the cap rather than present the table as a complete sweep; the $N=3$
coverage is three arms of nine, chosen as baseline, the predicted-clean case, and
the most aggressive failing case.
\textbf{The plant fires.}  Arm \texttt{atp} divides $\At_{\bM,\bI}$ by
$(w_1-w_2)$ for $|\bM|$ odd.  This is \emph{not} a point-product renormalisation
and it deliberately violates the unit hypothesis of Lemma~\ref{lem:unit}, so it
must produce $A$-type poles or the $\mathbf{0}$ column above is decoration.  It
produces $13$ of them at $N=2$, every one classified \texttt{A} with the correct
index pair $(1,2)$ and the correct factor $w_1-w_2$, on the \emph{same} solve and
the \emph{same} classifier that return $0$ for all eight admissible arms.  That
is the two-arm structure the $\mathbf{0}$ column needs: the instrument is
demonstrably able to see the thing it is reporting absent.

The \texttt{atp} plant probes one hypothesis only --- it introduces a \emph{pole},
so it tests ``$\Psi$ regular'' and nothing else.  The remaining hypotheses are
lesioned separately in V7, which is where the sharp controls are.

\paragraph{V7. Lesioning the theorem's own hypotheses (\texttt{lesion.py}).}
Theorem~\ref{thm:B} rests on exactly three things; each is lesioned once.  A
lesion that \emph{survives} is a stronger theorem for free; one that
\emph{breaks} shows the hypothesis is load-bearing rather than decorative.  All at
$N=2$, type $C$; baseline $20$ nonzero, $1$ non-Laurent, \texttt{A}-set empty.

\emph{L1 --- is $\rho^A=1$ needed, or would any constant do?}  Algebraic, no
computation: if the multiplier were a constant unit $c$, so that
$\Gkm_c=\{f:f(\bM)\equiv c\,f(s\bM)\}$, then
$(fg)(\bM)\equiv c^{2}(fg)(s\bM)$ --- weight $c^{2}\neq c$ unless $c=1$, so
$\Gkm_c$ is \emph{not} closed under multiplication.  $\rho^A=1$ is necessary, and
this is literally the $\rho^2-\rho$ mechanism of Remark~\ref{rem:mech}.

\emph{L2 --- the colength match.}  Multiply \emph{one} basis column $\phi_{\bI_0}$
by $(w_1-w_2)$.  This keeps every entry regular (so it does not touch
Lemma~\ref{lem:reg}), is \emph{label}-indexed, and raises
$\sum_\bI\ord\At_{\bI,\bI}$ from $2^{N-2}$ to $2^{N-2}+1$, making $\Lat$ a
colength-$1$ \emph{proper} sublattice of $\Gkm$.  Measured:
\begin{center}\small
\begin{tabular}{lrrl}
\toprule
$\bI_0$ & non-Laurent & \texttt{A}-type & \\\midrule
$\emptyset$ & $1$ & $0$ & structurally spared --- see below \\
$\{1\}$     & $2$ & $1$ & \textsc{fires} \\
$\{2\}$     & $4$ & $3$ & \textsc{fires} \\
$\{1,2\}$   & $6$ & $6$ & \textsc{fires} \\
\bottomrule
\end{tabular}
\end{center}
So the colength match is load-bearing: break it by one and $A$-type poles appear.
The $\bI_0=\emptyset$ row is spared for a nameable reason, not because the control
is blind: $\emptyset$ is minimal in Motegi's order, so $\At_{\emptyset,\lambda}=0$
for every $\lambda\neq\emptyset$ by (M3); hence $c^{\emptyset}_{\lambda\mu}$ is
either zero or carries a compensating factor of $(w_1-w_2)$ from the numerator.
Triangularity deletes the lesion exactly on the one column where it would have to
act alone.

\emph{L3 --- point-indexedness, and what the hypothesis really is.}  Two weights,
both units at every point, neither a point-indexed product:
\begin{center}\small
\begin{tabular}{llrrl}
\toprule
$\Psi$ & in $\Gkm^{\times}$? & non-Laurent & \texttt{A}-type & \\\midrule
$1+(w_1-w_2)$ on $\bM=\{1\}$, $1$ else & yes ($\equiv1$) & $5$ & $\mathbf{0}$ & \textsc{survives} \\
$2$ on $\bM=\{1\}$, $1$ else            & \textbf{no} ($2\not\equiv1$) & $7$ & $\mathbf{6}$ & \textsc{fires} \\
\bottomrule
\end{tabular}
\end{center}
This pair is the reason Theorem~\ref{thm:B} is stated for $\Gkm^{\times}$ rather
than for the point-product family: the first lesion survives, so point-indexedness
was never the hypothesis, and the second fires, so pointwise invertibility is not
enough either.  The two weights differ \emph{only} in whether the perturbation is
$\equiv1$ modulo the uniformiser.  Note also that the first lesion's five
non-Laurent constants are all \texttt{bdy} --- it does disturb the
\emph{boundary} divisors, as it must, since nothing here constrains those; it is
the $A$-type column that stays empty.

\paragraph{V8. Theorem~\ref{thm:A}'s \emph{conclusion}, tested directly
(\texttt{lattice\_equality.py}).}
V4 checks the two colengths, which is the \emph{proof's} arithmetic; this checks
the equality $\Lat_{ij}=\Gkm_{ij}$ itself.  The nontrivial inclusion is
$\Gkm\subseteq\Lat$: take a generator of $\Gkm_{ij}$ over $\mathcal{O}_{ij}$ ---
the $s_{ij}$-symmetric orbit sums $\mathbf{1}_{\bM}+\mathbf{1}_{s\bM}$, the
elements $t\cdot\mathbf{1}_{\bM}$, and $\mathbf{1}_{\bM}$ for $s_{ij}$-fixed
$\bM$ --- expand it in the $\phi$ basis by the \emph{same} triangular solve the
structure constants use, and require every coefficient to have
$\ord\ge0$.  At $N=2$: $4/4$ generators clean in type $C$ and $4/4$ in type $B$
(an independent corroboration of Lemma~\ref{lem:typeB}).

The control is $\mathbf{1}_{\bM}$ for a \emph{moved} $\bM$, which is \emph{not}
in $\Gkm$ ($1-0=1\notin t\mathcal{O}$); Theorem~\ref{thm:A} then forces some
coefficient to have $\ord=-1$.  It \textsc{fires}, $1/1$ in both types.  Without
it the clean column would only say the solve returns regular coefficients for
\emph{everything} I happened to feed it.  The $N=3$ arm was still running when
this was written.

%=====================================================================
\section{Gaps and questions}

\begin{rem}[the type-$B$ gap I expected to file, and did not]\label{gap:B}
I had this paragraph drafted as a Gap: Lemmas~\ref{lem:sym}, \ref{lem:reg}
and~\ref{lem:A1} are type-free, but Lemma~\ref{lem:colenL} was proved from the
type-$C$ root product \eqref{eq:diag}, and for type $B$ I had only the
\emph{measurement} of V4 ($\sum_\bI\ord=2^{N-2}$ and the per-label formula, at
$N=2,3$ only, because the engine has no type-$B$ diagonal root product).  Reading
L1307--1340 closed it: Lemma~\ref{lem:typeB}.  Theorems~\ref{thm:A}
and~\ref{thm:B} are therefore \emph{proved} for both types, and the type-$B$ rows
of V4 are now a check on a proof rather than its only support.  Q514, posed at
the start of this session for exactly this gap, is answered within it and
retired.
\end{rem}

\begin{gapthm}[$N\ge5$ is not measured]
The symbolic census is run at $N\le4$ (Lemma~A1) and $N\le3$ (structure
constants).  The \emph{proofs} are uniform in $N$, so this is not a gap in the
argument; I record it because the colength match $2^{N-2}=2^{N-2}$ is the one
step where an off-by-one in the orbit count would be invisible at small $N$ ---
at $N=2$ both sides are $1$, and $1$ is the value a count of almost anything can
take.  $N=4$ is the first value at which $2^{N-2}=4$ is distinguishable from the
alternatives $N-1=3$ and $2^{N-2}+1=5$, and both halves of the colength match are
measured there: Lemma~A1 at $768$ instances over $48$ orbits (V3), and
$\sum_\bI\ord\At_{\bI,\bI}=4$ for all six pairs with the per-label formula holding
on all $16$ labels (V4).  At $N=2$ both sides are $1$, and $1$ is the value a
count of almost anything can take.
\end{gapthm}

\begin{gapthm}[what this does \emph{not} say about $\Ah$ versus Koncki's class]
Corollary~\ref{cor:inv} is a statement about point-indexed \emph{diagonal}
renormalisations.  Rick's UID 824 result is that Koncki's single-entry
discrepancy at $N=2$ is \emph{unipotent off-diagonal in label space},
$\mathrm{Koncki}_\emptyset=\At_\emptyset+c\,\At_{\{1,2\}}$ with
$c=q(1-q)/(1-1/(w_1w_2))$ having a simple pole on the pair divisor $H_{12}$.
That lies outside the family here and outside the previous paper's Theorem~C.
The $A$-type direction cannot contribute to the Q510 identification at all,
precisely because it is blind to the whole admissible family.
\end{gapthm}

\paragraph{New questions.}
\textbf{Q513} --- does ``spans exactly the GKM lattice at all three root
families'' pin $\Ah$ up to a point-product unit?  Theorem~\ref{thm:A} supplies
the $A$-type lattice equality and the previous paper's Theorem~B supplies the
pair-divisor colength argument.  If the three equalities plus (M3) plus
\eqref{eq:diag} determine the basis up to $\Psi_\chi$ with $\chi$ a unit at every
divisor, then Q510 becomes a theorem about characterisations rather than a table
comparison.  Gate: the $A$-type direction imposes no condition, so count the
remaining freedom before claiming a characterisation.
\textbf{Q514} --- \emph{answered and retired within this session}
(Remark~\ref{gap:B}): \texttt{normalizationWB} (L1386) is indeed a $B_N$
positive-root product, its $A$-type factors are again
$\prod_{j\in\bI^{c}}\prod_{l:i_l>j}(1-w_{i_l}^{-1}w_j)$, and the whole
type-dependence of the diagonal is the boundary root $2t_k$ versus $t_k$.  See
Lemma~\ref{lem:typeB}.

%=====================================================================
\section*{Acknowledgement}

Rick's Lemma QC (UID 819) and his independent GKM check of $\Ah$ at $N=2$ (UID
824, run on \emph{his} instrument, which is what makes it corroboration rather
than my own evidence counted twice) set up the question this paper answers.  He
explicitly declined $N\ge3$; that was this session's target and
Theorem~\ref{thm:B} covers all $N$.

\end{document}
